% ===================================================================
%  TAVOLA N. 9 — La montagna e le terme. --- Il bosco e la caccia.
%  Traduction 2026 d'apres texto/io, controlee sur texto/fr et
%  texto/en. Consigne : voir texto/it/10-tavola-01.tex.
%
%  UN MOT DE CETTE PAGE FAIT LE TOUR DE L'EUROPE ET REVIENT CHEZ LUI
%  SANS SE RECONNAITRE.
%
%  L'alinea 2 parle des avalanches. Le tableau 9 allemand avait
%  releve que l'allemand des Alpes a exporte son vocabulaire de haute
%  montagne dans vingt langues — Gletscher, Lawine, Alpenstock,
%  Rucksack — et il avait note en passant que « Lawine » venait du
%  romanche « lavina ».
%
%  Or « lavina » est le latin « labina », de « labi », glisser, et
%  c'est l'italien valanga par une autre route. LE MOT EST
%  PARTI DU VERSANT SUD, EST MONTE CHEZ LES GERMANOPHONES, A PRIS
%  L'HABIT ALLEMAND, ET EST REPARTI DE LA DANS TOUTE L'EUROPE. Le
%  francais, l'anglais et le russe l'ont recu de l'allemand ; l'italien
%  ne l'a jamais emprunte parce qu'il ne l'avait jamais perdu.
%
%     UN MOT PEUT FAIRE LE TOUR DU MONDE ET REVENIR SANS QUE PERSONNE
%     NE VOIE QU'IL EST RENTRE. Le tableau 11 allemand avait trouve la
%     meme figure avec « Bank », parti d'Allemagne comme meuble et
%     revenu d'Italie comme metier. Ici c'est l'inverse, et le voyage
%     est le meme.
%
%  ET LA SECONDE SCENE DIT COMMENT L'ITALIEN NOMME LE BOIS, ET
%  POURQUOI IL LE NOMME MAL.
%
%  Le tableau 9 allemand avait donne la foret comme le quatrieme des
%  ilots ou l'allemand n'a rien emprunte — Wald, Buche, Birke, Eiche,
%  Hirsch, Fuchs, Dachs, vingt-trois mots, pas un emprunt. L'italien
%  a trois mots pour la meme chose, et DEUX SONT ETRANGERS :
%     selva   < silva, le seul qui soit latin ;
%     bosco   germanique, du lombard « busk » ;
%     foresta < « forestis silva », le bois qui est FORIS,
%                      DEHORS — hors du droit commun, reserve au roi.
%
%  Le mot que l'Europe entiere emploie — forest, Forst, forêt,
%  floresta — est donc un terme de chancellerie franque, et il ne
%  decrit pas un lieu mais une INTERDICTION. LES ALLEMANDS ONT NOMME
%  LE BOIS OU ILS VIVAIENT ; LES ITALIENS ONT NOMME LE BOIS DONT ON
%  LES TENAIT DEHORS.
%
%  La colonne ecrit « bosco » partout dans la seconde scene, parce que
%  c'est le mot ordinaire, et il se trouve que c'est le mot lombard.
%  Le seul mot latin des trois, « selva », ne sert plus qu'en poesie.
% ===================================================================

%%K t09-apar-1 apar
\VUsaut{4.0mm}
\VUtitre{15.0pt}{60}{TAVOLA N. 9}
\VUsaut{3.0mm}

%%K t09-tit-1 sub
\VUcentre{12.0pt}{0}{\textit{Prima scena.}}
\VUsaut{3.0mm}

%%K t09-tit-2 sub
\VUpk{11.4pt}{40}{La montagna. --- La città termale.}
\VUsaut{3.0mm}

%%K t09-c1-01-1 p
1. --- Sono a Pratz da una settimana. Non riesco a dire tutto lo stupore che ho provato
davanti alla meravigliosa \VUgras{catena} \textsuperscript{(1)} dei Pirenei, la cui
\VUgras{cresta} \textsuperscript{(2)} si staglia sul cielo azzurro come una gigantesca
conchiglia.

%%K t09-c1-02-1 p
2. --- Non avevo immaginato che le \VUgras{vette} \textsuperscript{(3)} dei Pirenei
potessero raggiungere un'\VUgras{altezza} simile, e sono rimasto stupito davanti ai
magnifici \VUgras{ghiacciai} \textsuperscript{(5)} e alle \VUgras{cime} innevate
\textsuperscript{(4)} da cui scendono \VUgras{valanghe} così spaventose.

%%K t09-tit-3 sub
\VUsaut{3.0mm}
\VUpk{10.2pt}{40}{Paesaggio di montagna.}
\VUsaut{3.0mm}

%%K t09-c1-03-1 p
3. --- Il giorno dopo il mio arrivo sono andato al vecchio \VUgras{vulcano}
\textsuperscript{(6)} spento vicino a Pratz. Sull'orlo nudo e brullo del
\VUgras{cratere} si vedono ancora chiaramente le tracce lasciate dalla \VUgras{lava}
che è colata, alcuni secoli fa, giù per il \VUgras{fianco della montagna}
\textsuperscript{(7)}. Su uno dei \VUgras{pendii} sono riuscito a trovare qualche
pezzo di quella lava.

%%K t09-c1-04-1 p
4. --- Pratz è al confine, e la \VUgras{fortezza} \textsuperscript{(8)} che sorveglia
il \VUgras{passo} \textsuperscript{(27)} fra la Francia e la Spagna somiglia molto a
quei vecchi castelli sulle rive del Reno che sembrano spiare la \VUgras{preda} dalla
cima della loro roccia.

%%K t09-c1-05-1 p
5. --- Un'\VUgras{aquila} \textsuperscript{(9)} enorme, che ha il nido non lontano dal
\VUgras{forte} \textsuperscript{(8)}, attira spesso la mia attenzione. Questo
\VUgras{rapace}, con il suo \VUgras{becco} potente, porta spesso ai suoi piccoli un
agnello o un capretto stretto negli \VUgras{artigli}. Le sue \VUgras{ali} sono così
grandi che può planare a lungo con il suo carico pesante.

%%K t09-tit-4 sub
\VUsaut{3.0mm}
\VUpk{10.2pt}{40}{L'escursionista.}
\VUsaut{3.0mm}

%%K t09-c1-06-1 p
6. --- La passeggiata più bella di qui è la gita alla \VUgras{cascata}
\textsuperscript{(10)} di Cau. L'\VUgras{acqua che cade} precipita in un vortice e poi
scompare come un grazioso \VUgras{ruscello} nella \VUgras{pineta}
\textsuperscript{(11)} a ovest della città. Abbiamo salito questo bosco fino
all'\VUgras{osservatorio meteorologico} \textsuperscript{(12)}, e dalla cima della
\VUgras{terrazza} abbiamo abbracciato tutto il \VUgras{panorama} della zona. Il
\VUgras{paesaggio} è superbo, e la \VUgras{natura} sembra aver versato su questo paese
tutti i tesori di cui disponeva.

%%K t09-c1-07-1 p
7. --- Per scendere abbiamo preso la \VUgras{funicolare} \textsuperscript{(13)} e ci
siamo fermati a una piccola \VUgras{stazione} vicino ai \VUgras{ruderi}
\textsuperscript{(14)} di un vecchio \VUgras{castello feudale} abitato un tempo dai
signori di Pratz.

%%K t09-c1-08-1 p
8. --- Povero castello! Che cosa penserà quando vede sotto di sé la graziosa
\VUgras{città termale} \textsuperscript{(15)}, diventata ormai una \VUgras{stazione
alla moda}?

%%K t09-c1-08-2 p
Il \VUgras{clima} è così piacevole e l'\VUgras{acqua minerale} così benefica che la
\VUgras{cura} che si fa al \VUgras{sanatorio} \textsuperscript{(17)} è un vero piacere.

%%K t09-tit-5 sub
\VUsaut{3.0mm}
\VUpk{10.2pt}{40}{Una graziosa città termale.}
\VUsaut{3.0mm}

%%K t09-c1-09-1 p
9. --- Si è fatto davvero di tutto per rendere gradevole il soggiorno. È stato
costruito un grande \VUgras{viadotto} \textsuperscript{(16)} per una ferrovia; il
\VUgras{parco} \textsuperscript{(18)} dello \VUgras{stabilimento}, dove si tengono i
\VUgras{concerti}, è disposto con molto gusto. Diverse graziose \VUgras{ville}
\textsuperscript{(19)} sono state costruite intorno al \VUgras{casinò}
\textsuperscript{(21)}, e una \VUgras{strada} \textsuperscript{(22)} tenuta molto bene
rende gli spostamenti estremamente facili.

%%K t09-c1-10-1 p
10. --- Un semplice \VUgras{terrapieno} \textsuperscript{(23)} separa la
\VUgras{strada} \textsuperscript{(20)} dalla \VUgras{valle} \textsuperscript{(25)} vera
e propria, in fondo alla quale c'è un grazioso \VUgras{laghetto}
\textsuperscript{(26)} che alimenta un limpido \VUgras{ruscello}
\textsuperscript{(24)}.

%%K t09-c1-11-1 p
11. --- Dall'altra parte della valle è in funzione una \VUgras{miniera di ferro}
\textsuperscript{(28)}. Più volte sono andato a vedere i \VUgras{minatori} che
scendevano nel \VUgras{pozzo}, e sono perfino entrato nella \VUgras{galleria} dove
passano la giornata a estrarre il \VUgras{minerale} che contiene il \VUgras{metallo}.

%%K t09-c1-12-1 p
12. --- Accanto alla miniera è stata costruita una grande \VUgras{ferriera} per
sfruttare subito le ricchezze che se ne cavano. L'\VUgras{altoforno}
\textsuperscript{(29)}, scaldato con il \VUgras{carbone} che qui abbonda, permette di
ottenere in fretta e a poco prezzo la \VUgras{ghisa}.

%%K t09-c1-13-1 p
13. --- Non lontano dalla miniera si scava una \VUgras{cava} \textsuperscript{(30)}.
Quello che il vecchio \VUgras{cavatore} stacca con il suo \VUgras{piccone} non è né
\VUgras{granito}, né \VUgras{porfido}, né \VUgras{arenaria}, ma semplice
\VUgras{pietra da costruzione} per fare le case.

%%K t09-c1-14-1 p
14. --- Uno degli ornamenti più belli di questo paese è l'\VUgras{abete}
\textsuperscript{(31)}. Dovunque --- in fondo a un \VUgras{burrone}
\textsuperscript{(32)}, sulla riva di un \VUgras{torrente} \textsuperscript{(33)},
accanto a un \VUgras{abisso} \textsuperscript{(34)} o a un precipizio --- i suoi aghi
sottili mettono la loro nota verde.

%%K t09-tit-6 sub
\VUsaut{3.0mm}
\VUpk{10.2pt}{40}{Camminare.}
\VUsaut{3.0mm}

%%K t09-c1-15-1 p
15. --- Approfitto delle vacanze per fare lunghe camminate. I \VUgras{sentieri}
\textsuperscript{(35)} che salgono a zigzag per la montagna permettono di fare diversi
\VUgras{chilometri}, e spesso diverse \VUgras{leghe}, senza accorgersi della distanza.

%%K t09-c1-15-2 p
\VUgras{Pietre miliari} \textsuperscript{(36)} e \VUgras{segnavia}
\textsuperscript{(37)} indicano i sentieri, sui quali si incontra spesso un giovane
\VUgras{montanaro} \textsuperscript{(38)} con un \VUgras{tascapane}
\textsuperscript{(40)} al fianco, che porta con sé una \VUgras{marmotta}
\textsuperscript{(39)}, oppure qualche \VUgras{zingaro} \textsuperscript{(41)} il cui
\VUgras{colbacco} \textsuperscript{(42)} fa uno strano effetto accanto al suo vecchio
\VUgras{cappuccio} \textsuperscript{(43)} strappato. Porta un \VUgras{orso}
\textsuperscript{(44)} ammaestrato alla \VUgras{catena}.

%%K t09-tit-7 sub
\VUpk{10.2pt}{40}{Alpinismo.}
\VUsaut{3.0mm}

%%K t09-c1-16-1 p
16. --- Quasi ogni giorno vado con una \VUgras{guida} \textsuperscript{(48)} a fare una
\VUgras{salita}, e sono molto fiero quando, con il mio \VUgras{sacco}
\textsuperscript{(47)} in spalla e il mio \VUgras{bastone ferrato} in mano, come un
vero \VUgras{turista} \textsuperscript{(46)}, attacco le \VUgras{rocce}
\textsuperscript{(45)}.

%%K t09-c1-17-1 p
17. --- Dalla cima di una montagna la gente della pianura non sembra più grande di
formiche; un \VUgras{gregge di pecore} \textsuperscript{(50)} che pascola in un
\VUgras{pascolo} \textsuperscript{(49)} sembra un giocattolo. Un \VUgras{pino}
\textsuperscript{(51)}, o perfino uno \VUgras{chalet di legno} \textsuperscript{(52)},
è poco più di un puntino sul verde.

%%K t09-c1-18-1 p
18. --- In queste gite incontriamo spesso sul \VUgras{sentiero}
\textsuperscript{(53)} un \VUgras{cacciatore di camosci} \textsuperscript{(54)} che
porta sulla spalla l'animale appena ucciso.

%%K t09-c1-19-1 p
19. --- Nel mio erbario ho messo una fronda di \VUgras{felce} \textsuperscript{(55)}
molto notevole e un grazioso rametto rosa di \VUgras{erica} \textsuperscript{(57)},
che ho colto all'imbocco di una piccola \VUgras{grotta} \textsuperscript{(56)}
nell'ultima camminata.

%%K t09-tit-8 sub
\VUsaut{3.0mm}
\VUcentre{12.0pt}{0}{\textit{Seconda scena.}}
\VUsaut{3.0mm}

%%K t09-tit-9 sub
\VUpk{11.4pt}{40}{Il bosco. --- La caccia.}
\VUsaut{3.0mm}

%%K t09-c2-01-1 p
1. --- Ieri ho passato una giornata magnifica nel bosco. La vita affaccendata degli
animali e degli \VUgras{insetti} \textsuperscript{(5)} che ci vivono mi ha affascinato.

%%K t09-c2-01-2 p
Il \VUgras{ragno} \textsuperscript{(1)} che tesse la sua \VUgras{tela}
\textsuperscript{(2)} fra due rami per prendere la \VUgras{mosca}
\textsuperscript{(3)} di passaggio, la \VUgras{lumaca} \textsuperscript{(4)} che si
trascina a fatica con il suo guscio, il cervo volante a caccia di preda, la formica
indaffarata accanto al \VUgras{formicaio} \textsuperscript{(6)} --- tutte queste
piccole creature andavano e venivano e lavoravano con un'energia straordinaria.

%%K t09-c2-02-1 p
2. --- Un \VUgras{serpente} \textsuperscript{(7)}, che all'inizio avevo preso per una
\VUgras{vipera} ma che era solo una \VUgras{biscia} innocua, stava arrotolato sotto un
tronco e ogni tanto tirava fuori la sua testolina sottile, che sembrava sempre pronta a
mordere. Davanti a me una grossa \VUgras{lumaca senza guscio} \textsuperscript{(8)}
strisciava pesantemente dopo aver divorato una di quelle saporite piccole
\VUgras{fragole di bosco} \textsuperscript{(11)} che qui crescono in abbondanza.

%%K t09-c2-03-1 p
3. --- Il terreno era tappezzato di \VUgras{fiori di bosco}:
\VUgras{primule} \textsuperscript{(9)}, \VUgras{pervinche} \textsuperscript{(10)} e
molte altre piante che non sapevo nominare fiorivano fra i \VUgras{ciuffi d'erba}
\textsuperscript{(12)}.

%%K t09-c2-04-1 p
4. --- In questa zona il \VUgras{tartufo} è quasi comune quanto il \VUgras{fungo}
\textsuperscript{(14)}, e un \VUgras{maialino} \textsuperscript{(13)} di un guardiano
cercava avidamente se ne trovava.

%%K t09-tit-10 sub
\VUsaut{3.0mm}
\VUpk{10.2pt}{40}{Il bracconaggio.}
\VUsaut{3.0mm}

%%K t09-c2-05-1 p
5. --- Ho visto la \VUgras{fine} di una scena che mi ha fatto molto ridere. Grazie al
fiuto del suo \VUgras{bassotto} \textsuperscript{(15)}, il \VUgras{guardiano}
\textsuperscript{(16)} aveva appena sorpreso un \VUgras{bracconiere}
\textsuperscript{(22)} nel momento in cui l'uomo si preparava a portar via la
\VUgras{selvaggina} \textsuperscript{(17)} che aveva ucciso. Preferendo lasciare la
preda piuttosto che farsi arrestare, il bracconiere era scappato, e \VUgras{lepre}
\textsuperscript{(18)}, \VUgras{cerva} \textsuperscript{(19)},
\VUgras{capriolo} \textsuperscript{(20)} e \VUgras{cinghiale}
\textsuperscript{(21)} erano rimasti sull'erba, con gran soddisfazione del guardiano.

%%K t09-c2-06-1 p
6. --- La varietà degli alberi che il bosco contiene è sorprendente. I \VUgras{faggi}
\textsuperscript{(23)} e le \VUgras{betulle} \textsuperscript{(26)}, i \VUgras{pini},
che danno la \VUgras{resina} \textsuperscript{(27)}, le \VUgras{querce}
\textsuperscript{(29)} e molti altri ancora mescolano i loro rami.

%%K t09-c2-06-2 p
L'\VUgras{edera} \textsuperscript{(24)} che si aggrappa ai tronchi degli alberi alti dà
alla \VUgras{fustaia} \textsuperscript{(25)} un verde brillante, che si accorda
piacevolmente con il tono più pallido e più morbido del \VUgras{musco}
\textsuperscript{(31)} in cui il \VUgras{picchio verde} \textsuperscript{(30)} cerca
gli insetti.

%%K t09-tit-11 sub
\VUsaut{3.0mm}
\VUpk{10.2pt}{40}{Paesaggio di bosco.}
\VUsaut{3.0mm}

%%K t09-c2-07-1 p
7. --- Le vecchie querce mi hanno incantato. Le loro grosse \VUgras{radici}
\textsuperscript{(32)} contorte, mezze fuori dal terreno, danno loro un aspetto
splendido di forza e di vigore, e mi chiedo come il \VUgras{proprietario}
\textsuperscript{(35)} possa decidersi a farle abbattere e a consegnarle alla
\VUgras{sega} \textsuperscript{(34)} del \VUgras{boscaiolo} \textsuperscript{(33)}. I
poveri \VUgras{tronchi} \textsuperscript{(36)}, spogliati della loro \VUgras{corteccia}
\textsuperscript{(37)} o spaccati dall'\VUgras{ascia} \textsuperscript{(38)} del
\VUgras{giornaliero} \textsuperscript{(39)}, mi fanno pena.

%%K t09-c2-08-1 p
8. --- Lì vicino un vecchio \VUgras{carbonaio} \textsuperscript{(41)} aveva alzato con
la sua \VUgras{pala} una \VUgras{carbonaia} \textsuperscript{(42)}, e una vecchia
portava via una \VUgras{fascina} \textsuperscript{(40)} di \VUgras{legna secca} fatta
con i rami degli alberi abbattuti.

%%K t09-tit-12 sub
\VUsaut{3.0mm}
\VUpk{10.2pt}{40}{La caccia.}
\VUsaut{3.0mm}

%%K t09-c2-09-1 p
9. --- Volevo riposarmi un po' alla \VUgras{casa del guardiano} \textsuperscript{(46)},
ma quando sono arrivato al piccolo \VUgras{stagno} \textsuperscript{(43)}, su cui
nuotava un bel \VUgras{cigno} \textsuperscript{(45)}, mi sono accorto che una
\VUgras{piena} \textsuperscript{(44)} recente aveva trasformato il sentiero in una
vera \VUgras{palude}. Ho dovuto fare un giro, e mentre passavo accanto al
\VUgras{cedro} \textsuperscript{(49)}, ai \VUgras{pioppi} \textsuperscript{(48)} e
all'\VUgras{olmo} \textsuperscript{(47)} che stanno ai margini del bosco, ho visto
uscire da un \VUgras{fitto} \textsuperscript{(54)} un magnifico \VUgras{cervo}
\textsuperscript{(50)}, inseguito da una \VUgras{muta} \textsuperscript{(51)}
implacabile spronata dal \VUgras{corno} \textsuperscript{(53)} dei \VUgras{cacciatori
a cavallo} \textsuperscript{(52)}. La povera bestia era allo stremo. È caduta morente
a pochi passi da me.

%%K t09-c2-10-1 p
10. --- Il figlio del guardiano, che il rumore della \VUgras{caccia} aveva fatto
uscire, è venuto fuori di casa, e siamo rientrati in due nel bosco. Aveva visto una
\VUgras{gazza} \textsuperscript{(56)} sul ramo di una vecchia quercia. Pensava che il
nido fosse nascosto nel \VUgras{fogliame} \textsuperscript{(58)} e voleva prenderlo.

%%K t09-c2-10-2 p
Agile come uno \VUgras{scoiattolo} \textsuperscript{(61)}, si è arrampicato di
\VUgras{ramo} \textsuperscript{(55)} in ramo, ma non ha trovato niente ed è riuscito
solo a far scappare una vecchia \VUgras{cornacchia} \textsuperscript{(60)} posata su un
\VUgras{ramoscello} \textsuperscript{(57)}.
\VUsaut{4.0mm}
\VUfilet{22mm}
